This technical report benchmarked on two Speaker Embedding models: ECAPA-TDNN and RawNet3 under TidyVoice Vietnamese multi-dialectal (ViMD) corpora. On TidyVoice, ECAPA-TDNN demonstrated better cross-lingual generalization (3.912\% EER, 0.3586 minDCF, and 71.84\% TAR@0.1\%). However, when adapted to regional Vietnamese speech on ViMD, Raw-Net3 proved more efficacy at disentangling dialect from identity (3.071\% EER, 0.2262 minDCF, and 84.45\% TAR), which also comes with the benefit of smaller parameter footprint. Thus, RawNet3 was selected as the biometric engine for our Vietnamese target application.

Translating these empirical findings into practice, we developed \emph{Who-Speak-AI}, an end-to-end voice assistant integrating RawNet3 into a modular pipeline featuring Whisper ASR, LLM dialogue reasoning, and neural TTS. The system couples RawNet3 embeddings with the CKKS homomorphic encryption scheme, allowing server-side cosine matching to execute strictly over encrypted ciphertexts. In our experiments, this approach maintains high accuracy while protecting user privacy.

Nevertheless, several limitations remain: \textbf{(1)} only two models were benchmarked, as WavLM + MHFA failed during training and testing; \textbf{(2)} fine-tuning was constrained to a single seed (42) and 3 epochs without exhaustive tuning; \textbf{(3)} thresholds were calibrated on static splits rather than live streams; and \textbf{(4)} the system with homomorphic operations currently CPU-bound. Future work will explore anti-spoofing and hardware acceleration for encrypted matching.